%chapters/ce1-2025-formulae.tex
% Labels added so every \eqref{} used in ce1-2025-problems.tex resolves.
% Lines marked  %% NEW  are additions/changes made for labelling.

\begin{enumerate}[label=\thesubsection.\arabic*, ref=\thesubsection.\theenumi,itemsep=0pt]

\item Intersection of Conics and Common Chord
\begin{itemize}
\item Conic in Vector Form:
Every second-degree curve $ax^2 + 2bxy + cy^2 + 2dx + 2ey + f = 0$ can be written in matrix-vector form as:
\begin{align}
\vec{x}^\top \vec{V} \vec{x} + 2\vec{u}^\top \vec{x} + f = 0, \quad \vec{V} = \myvec{a & b \\ b & c} = \vec{V}^\top, \quad \vec{u} = \myvec{d \\ e}
\label{eq:conic} %% NEW
\end{align}
The conic is classified by the determinant of $\vec{V}$:
\begin{align} %% NEW (was inline)
\abs{\vec{V}} = \mydet{a & b \\ b & c} = ac - b^2 \begin{cases} = 0, & \text{parabola} \\ > 0, & \text{ellipse} \\ < 0, & \text{hyperbola} \end{cases}
\label{eq:class} %% NEW
\end{align}
For a circle with centre $\vec{c}$ and radius $r$, $\norm{\vec{x} - \vec{c}}^2 = r^2$ gives $\vec{V} = \vec{I}$, $\vec{u} = -\vec{c}$, and $f = \norm{\vec{c}}^2 - r^2$.

\item Line-Conic Intersection:
The points of intersection of a line $L: \vec{x} = \vec{h} + \kappa \vec{m}$ ($\kappa \in \mathbb{R}$) with a conic section $\vec{x}^\top \vec{V} \vec{x} + 2\vec{u}^\top \vec{x} + f = 0$ are given by:
\begin{align}
\vec{x}_i = \vec{h} + \kappa_i \vec{m}
\end{align}
where the parameter values $\kappa_i$ are:
\begin{align}
\kappa_i = \frac{-\vec{m}^\top \brak{\vec{V}\vec{h} + \vec{u}} \pm \sqrt{\sbrak{\vec{m}^\top \brak{\vec{V}\vec{h} + \vec{u}}}^2 - g\brak{\vec{h}}\brak{\vec{m}^\top \vec{V} \vec{m}}}}{\vec{m}^\top \vec{V} \vec{m}}
\label{eq:kappa} %% NEW
\end{align}
with
\begin{align} %% NEW (was inline)
g\brak{\vec{h}} = \vec{h}^\top \vec{V} \vec{h} + 2\vec{u}^\top \vec{h} + f
\label{eq:gh} %% NEW
\end{align}
The line meets the conic at two, one, or no real points according as the discriminant
\begin{align} %% NEW (was inline; now displayed so it can be referenced)
\Delta = \sbrak{\vec{m}^\top \brak{\vec{V}\vec{h} + \vec{u}}}^2 - g\brak{\vec{h}}\brak{\vec{m}^\top \vec{V} \vec{m}}
\label{eq:disc} %% NEW
\end{align}
is positive, zero, or negative.

\item Derivation of the Intersection Parameter:
Substituting $\vec{x} = \vec{h} + \kappa \vec{m}$ into the conic equation and using symmetry $\vec{V} = \vec{V}^\top$:
\begin{align}
\brak{\vec{h} + \kappa \vec{m}}^\top \vec{V} \brak{\vec{h} + \kappa \vec{m}} + 2\vec{u}^\top \brak{\vec{h} + \kappa \vec{m}} + f &= 0 \\
\kappa^2 \brak{\vec{m}^\top \vec{V} \vec{m}} + 2\kappa \vec{m}^\top \brak{\vec{V}\vec{h} + \vec{u}} + g\brak{\vec{h}} &= 0
\label{eq:quadk} %% NEW
\end{align}
When $\Delta = 0$, the two roots coincide and the line is tangent to the conic.

\item Pencil of Conics:
The intersection of two conics with parameters $\brak{\vec{V}_1, \vec{u}_1, f_1}$ and $\brak{\vec{V}_2, \vec{u}_2, f_2}$ lies on every member of the pencil:
\begin{align}
\vec{x}^\top \brak{\vec{V}_1 + \mu \vec{V}_2}\vec{x} + 2\brak{\vec{u}_1 + \mu \vec{u}_2}^\top \vec{x} + \brak{f_1 + \mu f_2} = 0
\label{eq:pencil} %% NEW
\end{align}

\item Degenerate Conic and Common Chord:
The pencil represents a pair of straight lines (a degenerate conic) if and only if:
\begin{align}
\mydet{\vec{V}_1 + \mu \vec{V}_2 & \vec{u}_1 + \mu \vec{u}_2 \\ \brak{\vec{u}_1 + \mu \vec{u}_2}^\top & f_1 + \mu f_2} = 0
\label{eq:degen} %% NEW
\end{align}
If $\vec{V}_1 = \vec{V}_2$ (e.g., two circles), choosing $\mu = -1$ removes the second-degree terms, leaving the common chord (radical axis):
\begin{align}
2\brak{\vec{u}_1 - \vec{u}_2}^\top \vec{x} + \brak{f_1 - f_2} = 0
\label{eq:chord} %% NEW
\end{align}
\end{itemize}

\item Line Intersection and Angle Between Vectors
\begin{itemize}
\item Forms of a Line:
A line through point $\vec{A}$ with direction vector $\vec{m}$ can be written in parametric form or normal form using normal vector $\vec{n}$ ($\vec{n}^\top \vec{m} = 0$):
\begin{align}
\vec{x} = \vec{A} + \kappa \vec{m}, \quad \vec{n}^\top \brak{\vec{x} - \vec{A}} = 0
\label{eq:lineform} %% NEW
\end{align}
In 2D, if $\vec{m} = \myvec{m_1 \\ m_2}$, then
\begin{align} %% NEW (was inline)
\vec{n} = \myvec{-m_2 \\ m_1}
\label{eq:normal} %% NEW
\end{align}

\item Line Intersection and Solvability:
The lines $\vec{x} = \vec{A} + \kappa_1 \vec{m}_1$ and $\vec{x} = \vec{B} + \kappa_2 \vec{m}_2$ meet where:
\begin{align}
\myvec{\vec{m}_1 & -\vec{m}_2} \myvec{\kappa_1 \\ \kappa_2} = \vec{B} - \vec{A}
\label{eq:lineint} %% NEW
\end{align}
This system has a unique solution if and only if $\mydet{\vec{m}_1 & \vec{m}_2} \ne 0$. If the determinant is zero, the lines are parallel or coincident.

\item Angle Between Vectors:
The angle $\theta$ between two vectors $\vec{a}$ and $\vec{b}$ is given by:
\begin{align}
\cos\theta = \frac{\vec{a}^\top \vec{b}}{\norm{\vec{a}}\norm{\vec{b}}}
\label{eq:angle} %% NEW
\end{align}
Two vectors are orthogonal if and only if $\vec{a}^\top \vec{b} = 0$.

\item Circle Representation and Chords:
A point on a circle of radius $r$ centered at the origin is $\vec{x} = r \myvec{\cos\alpha \\ \sin\alpha}$ with $\vec{e}_\alpha \triangleq \myvec{\cos\alpha \\ \sin\alpha}$. The chord length between $\vec{x}_\alpha = r\vec{e}_\alpha$ and $\vec{x}_\beta = r\vec{e}_\beta$ satisfies:
\begin{align}
\vec{e}_\alpha^\top \vec{e}_\beta &= \cos\brak{\alpha - \beta} \label{eq:cosdiff} \\ %% NEW (line added + label)
\norm{\vec{x}_\alpha - \vec{x}_\beta}^2 &= 2r^2 \brak{1 - \vec{e}_\alpha^\top \vec{e}_\beta} = 4r^2 \sin^2 \frac{\alpha - \beta}{2}
\end{align}
\end{itemize}

\item System of Linear Equations and Rouch\'{e}--Capelli Theorem
\begin{itemize}
\item Consistency and Solvability:
A system
\begin{align} %% NEW (was inline; now displayed so it can be referenced)
\vec{A}\vec{x} = \vec{b}, \quad \vec{A} \in \mathbb{R}^{m \times n}
\label{eq:axb} %% NEW
\end{align}
is consistent if and only if
\begin{align}
\text{rank}\brak{\vec{A}} = \text{rank}\brak{{\vec{A} \mid \vec{b}}} = r
\label{eq:rc} %% NEW
\end{align}
If $r = n$, the solution is unique; if $r < n$, there exist infinitely many solutions with $n - r$ free variables. If $\vec{A} \in \mathbb{R}^{n \times \brak{n-1}}$, the system is consistent only if
\begin{align} %% NEW (was inline)
\mydet{\vec{A} & \vec{b}} = 0
\label{eq:augdet} %% NEW
\end{align}

\item Homogeneous System and Zero Norm Property:
For $\vec{A} \in \mathbb{R}^{2 \times 3}$ with rank 2, the solution of $\vec{A}\vec{x} = \vec{0}$ is given by the cross product of its rows:
\begin{align}
\vec{x} = \kappa \brak{\vec{a}_1 \times \vec{a}_2}
\label{eq:cross} %% NEW
\end{align}
For any vector $\vec{v} \in \mathbb{R}^n$,
\begin{align} %% NEW (was inline)
\norm{\vec{v}}^2 = \vec{v}^\top \vec{v} = \sum_i v_i^2 = 0 \iff \vec{v} = \vec{0}
\label{eq:zeronorm} %% NEW
\end{align}
\end{itemize}

\item Determinants and Cofactor Expansion %% NEW (whole item)
\begin{itemize}
\item Expansion Along a Row:
For $\vec{A} \in \mathbb{R}^{n \times n}$, expanding along row $i$ with minor $M_{ij}$ (the determinant after deleting row $i$ and column $j$):
\begin{align}
\mydet{\vec{A}} = \sum_{j=1}^n \brak{-1}^{i+j} a_{ij} M_{ij}
\label{eq:detexp} %% NEW
\end{align}
For a $2 \times 2$ matrix:
\begin{align}
\mydet{a & b \\ c & d} = ad - bc
\label{eq:det2} %% NEW
\end{align}
\end{itemize}

\item Relation Between Eigenvalues and Trace of a Matrix
\begin{itemize}
\item Characteristic Polynomial and Trace Relation:
The characteristic polynomial of $\vec{A} \in \mathbb{R}^{n \times n}$ is
\begin{align} %% NEW (was inline)
P\brak{\lambda} = \mydet{\lambda \vec{I} - \vec{A}}
\label{eq:charpoly} %% NEW
\end{align}
The sum of all eigenvalues equals the trace of the matrix:
\begin{align}
\sum_{i=1}^n \lambda_i = \sum_{i=1}^n a_{ii} = \operatorname{tr}\brak{\vec{A}}
\label{eq:trace} %% NEW
\end{align}
A matrix is symmetric if $\vec{A}^\top = \vec{A}$ and skew-symmetric if $\vec{A}^\top = -\vec{A}$.
\end{itemize}

\item Cholesky Decomposition
\begin{itemize}
\item Factorization:
A symmetric positive definite matrix $\vec{A}$ can be factorized as $\vec{A} = \vec{L}\vec{L}^\top$ \label{eq:chol1}  where $\vec{L}$ is lower triangular.
\item $2 \times 2$ Decomposition Form:
\begin{align}
\myvec{a_{11} & a_{12} \\ a_{12} & a_{22}} = \myvec{l_{11} & 0 \\ l_{21} & l_{22}} \myvec{l_{11} & l_{21} \\ 0 & l_{22}} = \myvec{l_{11}^2 & l_{11}l_{21} \\ l_{11}l_{21} & l_{21}^2 + l_{22}^2}
\label{eq:chol2} %% NEW
\end{align}
\end{itemize}

\item Bayes' Theorem in Vector Form
\begin{itemize}
\item Vector Formulation:
Given prior vector $\vec{\pi} = \myvec{\Pr\brak{B_1} & \dots & \Pr\brak{B_n}}^\top$ and likelihood vector $\vec{l} = \myvec{\Pr\brak{E \mid B_1} & \dots & \Pr\brak{E \mid B_n}}^\top$:
\begin{align}
\Pr\brak{E} &= \vec{\pi}^\top \vec{l} \label{eq:total} \\ %% NEW
\Pr\brak{B_i \mid E} &= \frac{\pi_i l_i}{\vec{\pi}^\top \vec{l}} \label{eq:bayes} %% NEW
\end{align}
\end{itemize}

\item Minimization via Completing the Square
\begin{itemize}
\item Inequality and Equality Condition:
For real $a$ and $b$:
\begin{align}
a b - b^2 = \frac{a^2}{4} - \brak{b - \frac{a}{2}}^2 \le \frac{a^2}{4}
\label{eq:cts} %% NEW
\end{align}
with equality holding if and only if $b = \frac{a}{2}$.
\item Zero Integral Property:
If $g\brak{x} \ge 0$ is continuous on $\sbrak{a,b}$ and $\int_{a}^{b} g\brak{x} dx = 0$, then $g\brak{x} = 0$ on $\sbrak{a,b}$.
\end{itemize}

\item Laplace Transform and State-Space Form
\begin{itemize}
\item Transform Operational Properties:
\begin{align}
y'\brak{t} &\system{L} s Y\brak{s} - y\brak{0} \label{eq:lap1} \\ %% NEW
y''\brak{t} &\system{L} s^2 Y\brak{s} - s y\brak{0} - y'\brak{0} \label{eq:lap2} \\ %% NEW
e^{at} u\brak{t} &\system{L} \frac{1}{s - a} \label{eq:lapexp} %% NEW
\end{align}
\item State-Space Representation:
Writing $\vec{z} = \myvec{y \\ y'}^\top$, the ODE $y'' + p y' + q y = 0$ becomes $\vec{z}' = \vec{A}\vec{z}$ where:
\begin{align}
\vec{A} = \myvec{0 & 1 \\ -q & -p}
\label{eq:ss} %% NEW
\end{align}
The exponents of the solution are the eigenvalues of $\vec{A}$.
\end{itemize}

\item Newton-Raphson Method
\begin{itemize}
\item First-Order Recurrence:
For a differentiable function $f\brak{x}$, the root is iteratively approximated by:
\begin{align}
x_{n+1} = x_n - \frac{f\brak{x_n}}{f'\brak{x_n}}
\label{eq:nr} %% NEW
\end{align}
\end{itemize}

\item Polynomial Interpolation (Vandermonde Matrix)
\begin{itemize}
\item Matrix Form:
The polynomial $P_2\brak{x} = c_0 + c_1 x + c_2 x^2$ through points $\brak{x_i, y_i}$ for $i=1,2,3$ satisfies:
\begin{align}
\myvec{1 & x_1 & x_1^2 \\ 1 & x_2 & x_2^2 \\ 1 & x_3 & x_3^2} \myvec{c_0 \\ c_1 \\ c_2} = \myvec{y_1 \\ y_2 \\ y_3}
\label{eq:vander} %% NEW
\end{align}
\end{itemize}

\item Equilibrium of Rigid Bodies in Vector Form
\begin{itemize}
\item Force and Moment Equilibrium:
For a body in equilibrium under forces $\vec{F}_i$ acting at points $\vec{r}_i$:
\begin{align}
\sum_i \vec{F}_i = \vec{0}, \quad \sum_i \brak{\vec{r}_i - \vec{O}} \times \vec{F}_i = \vec{0}
\label{eq:equil} %% NEW
\end{align}
where in the 2D plane, $\vec{r} \times \vec{F} = r_x F_y - r_y F_x$.
\item Bending Moment and Friction:
The bending moment at section $P$ is
\begin{align} %% NEW (was inline)
M\brak{P} = \sum_{\text{one side}} \brak{\vec{r}_i - \vec{P}} \times \vec{F}_i
\label{eq:bm} %% NEW
\end{align}
Limiting friction force is $F = \mu N$.
\end{itemize}

\item Direct Stiffness Method
\begin{itemize}
\item Member and Global Stiffness Formulation:
For an axial member of length $L$, area $A$, modulus $E$, and direction vector $\vec{m}$, the nodal stiffness is
\begin{align} %% NEW (was inline)
\vec{k} = \frac{AE}{L}\vec{m}\vec{m}^\top
\label{eq:kbar} %% NEW
\end{align}
and the member stiffness matrix is $\myvec{\vec{k} & -\vec{k} \\ -\vec{k} & \vec{k}}$. The global system is
\begin{align} %% NEW (was inline)
\vec{K}\vec{u} = \vec{F}
\label{eq:kuf} %% NEW
\end{align}
\item Thermal Load Equivalence:
A temperature rise $\Delta T$ with expansion coefficient $\alpha$ produces equivalent nodal loads of $AE\alpha\Delta T$ at member ends directed outward.
\end{itemize}
\item Free Thermal Expansion
\begin{itemize}
\item Thermal Strain and Elongation:
A temperature change $\Delta T$ in an unrestrained member of length $L$ with coefficient of thermal expansion $\alpha$ produces thermal strain $\epsilon_{th}$ and free thermal elongation $\Delta L$:
\begin{align}
\Delta L = L \alpha \Delta T
\label{eq:free}
\end{align}
\end{itemize}

\item Cellular Update Rule in Matrix Form
\begin{itemize}
\item Matrix Formulation:
For $N$ bulbs in a line with state vector $\vec{x}_k \in \{0, 1\}^N$, the number of ON neighbours is given by $\vec{s}_k = \vec{A}\vec{x}_k$ with:
\begin{align}
a_{ij} = \begin{cases} 1, & \abs{i - j} = 1 \\ 0, & \text{otherwise} \end{cases}
\label{eq:adj} %% NEW
\end{align}
\end{itemize}



\end{enumerate}
